% Options for packages loaded elsewhere
\PassOptionsToPackage{unicode}{hyperref}
\PassOptionsToPackage{hyphens}{url}
\PassOptionsToPackage{dvipsnames,svgnames,x11names}{xcolor}
%
\documentclass[
  10pt,
  a4paper,
]{article}
\usepackage{amsmath,amssymb}
\usepackage{iftex}
\ifPDFTeX
  \usepackage[T1]{fontenc}
  \usepackage[utf8]{inputenc}
  \usepackage{textcomp} % provide euro and other symbols
\else % if luatex or xetex
  \usepackage{unicode-math} % this also loads fontspec
  \defaultfontfeatures{Scale=MatchLowercase}
  \defaultfontfeatures[\rmfamily]{Ligatures=TeX,Scale=1}
\fi
\usepackage{lmodern}
\ifPDFTeX\else
  % xetex/luatex font selection
  \setmainfont[]{Noto Serif CJK JP}
  \setsansfont[]{Noto Sans CJK JP}
  \setmonofont[]{Noto Sans Mono CJK JP}
  \setmathfont[]{Latin Modern Math}
\fi
% Use upquote if available, for straight quotes in verbatim environments
\IfFileExists{upquote.sty}{\usepackage{upquote}}{}
\IfFileExists{microtype.sty}{% use microtype if available
  \usepackage[]{microtype}
  \UseMicrotypeSet[protrusion]{basicmath} % disable protrusion for tt fonts
}{}
\makeatletter
\@ifundefined{KOMAClassName}{% if non-KOMA class
  \IfFileExists{parskip.sty}{%
    \usepackage{parskip}
  }{% else
    \setlength{\parindent}{0pt}
    \setlength{\parskip}{6pt plus 2pt minus 1pt}}
}{% if KOMA class
  \KOMAoptions{parskip=half}}
\makeatother
\usepackage{xcolor}
\usepackage[margin=20mm]{geometry}
\usepackage{longtable,booktabs,array}
\usepackage{calc} % for calculating minipage widths
% Correct order of tables after \paragraph or \subparagraph
\usepackage{etoolbox}
\makeatletter
\patchcmd\longtable{\par}{\if@noskipsec\mbox{}\fi\par}{}{}
\makeatother
% Allow footnotes in longtable head/foot
\IfFileExists{footnotehyper.sty}{\usepackage{footnotehyper}}{\usepackage{footnote}}
\makesavenoteenv{longtable}
\setlength{\emergencystretch}{3em} % prevent overfull lines
\providecommand{\tightlist}{%
  \setlength{\itemsep}{0pt}\setlength{\parskip}{0pt}}
\setcounter{secnumdepth}{-\maxdimen} % remove section numbering
% XeLaTeX 用プリアンブル（pandoc の header-includes として読み込む）
% 日本語専用パッケージ（xeCJK / luatexja）なしで組めるよう、
% 本文フォントに和欧両対応の Noto Serif CJK JP を使い、
% 和文の行分割は XeTeX 組み込みのロケール改行に任せる。
\XeTeXlinebreaklocale "ja"
\XeTeXlinebreakskip=0pt plus 0.5pt
\XeTeXlinebreakpenalty=0
\usepackage{amsmath,amssymb}
\setlength{\parindent}{1em}
\setlength{\parskip}{0.4\baselineskip}
\renewcommand{\arraystretch}{1.25}
\usepackage{etoolbox}
\AtBeginEnvironment{longtable}{\small}

% Noto Serif CJK JP に無い記号（上付き・下付きの添字文字、ℝ・ℕ、⟨ ⟩ など）用の補助フォント。
% 該当箇所は latex/fallback-glyphs.lua が {\fallbackfont ...} で包む。
\newfontfamily\fallbackfont{DejaVu Serif}
\ifLuaTeX
  \usepackage{selnolig}  % disable illegal ligatures
\fi
\IfFileExists{bookmark.sty}{\usepackage{bookmark}}{\usepackage{hyperref}}
\IfFileExists{xurl.sty}{\usepackage{xurl}}{} % add URL line breaks if available
\urlstyle{same}
\hypersetup{
  colorlinks=true,
  linkcolor={blue},
  filecolor={Maroon},
  citecolor={Blue},
  urlcolor={blue},
  pdfcreator={LaTeX via pandoc}}

\author{}
\date{}

\begin{document}

\hypertarget{ux5faeux5206ux5e7eux4f55ux5b66-ux4eacux5927ux7269ux7406ux5f62ux5f0fux30c6ux30adux30b9ux30c8}{%
\section{微分幾何学
京大物理形式テキスト}\label{ux5faeux5206ux5e7eux4f55ux5b66-ux4eacux5927ux7269ux7406ux5f62ux5f0fux30c6ux30adux30b9ux30c8}}

作成日：2026-11-09

\hypertarget{ux4f7fux3044ux65b9}{%
\subsection{使い方}\label{ux4f7fux3044ux65b9}}

微分幾何学の基礎を6テーマに分け、京大物理と同じ「誘導つきの長文＋空欄補充＋選択＋問」形式で1テーマ＝大問1題にまとめた。前半3題で
{\fallbackfont ℝ}³
内の曲線・曲面の古典論を、後半3題で抽象的な多様体と微分形式を扱い、「曲面の内側だけで測れる量」から多様体論へ橋渡しする流れになっている。各大問は30〜40分が目安。

\begin{itemize}
\tightlist
\item
  【ア】などの空欄には式・数値・語句を入れる。前の空欄の答えを後で使うことがある。
\item
  \{ ケ \} などの選択欄には①〜③から1つ選ぶ。
\item
  「問」は導出・証明・計算。物理の記述問題と同じく、論理を短く書く。
\item
  記号は京大物理にならい、問題Ⅰ・Ⅳはカタカナ、問題Ⅱ・Ⅴはイロハ、問題Ⅲ・Ⅵはひらがな。
\item
  解答と解説は末尾にまとめた。
\end{itemize}

\begin{longtable}[]{@{}
  >{\raggedright\arraybackslash}p{(\columnwidth - 4\tabcolsep) * \real{0.3333}}
  >{\raggedright\arraybackslash}p{(\columnwidth - 4\tabcolsep) * \real{0.3333}}
  >{\raggedright\arraybackslash}p{(\columnwidth - 4\tabcolsep) * \real{0.3333}}@{}}
\toprule\noalign{}
\begin{minipage}[b]{\linewidth}\raggedright
大問
\end{minipage} & \begin{minipage}[b]{\linewidth}\raggedright
テーマ
\end{minipage} & \begin{minipage}[b]{\linewidth}\raggedright
到達目標
\end{minipage} \\
\midrule\noalign{}
\endhead
\bottomrule\noalign{}
\endlastfoot
Ⅰ & 曲線の幾何 & 弧長・曲率・捻率、フルネ＝セレの公式を計算できる \\
Ⅱ & 曲面と第一基本形式 & 長さ・面積・等長性を E, F, G で表せる \\
Ⅲ & 第二基本形式とガウス曲率 &
主曲率・ガウス曲率を計算し、曲面の形を分類できる \\
Ⅳ & 測地線・驚異の定理・ガウス＝ボネ &
曲率が内在的量であることと、その大域的帰結を理解する \\
Ⅴ & 可微分多様体と接空間 &
チャート・座標変換・接ベクトル（導分）を扱える \\
Ⅵ & 微分形式とストークスの定理 &
外微分、ストークスの定理、閉形式と完全形式の違い \\
\end{longtable}

共通の記号：曲面 p(u, v) の偏微分を p{\fallbackfont ᵤ},
p{\fallbackfont ᵥ}, p{\fallbackfont ᵤᵤ}, p{\fallbackfont ᵤᵥ},
p{\fallbackfont ᵥᵥ} と書き、「·」は {\fallbackfont ℝ}³
の内積、「×」は外積を表す。添字の多い式は数式ブロックで示す。

\hypertarget{ux554fux984cux2170-ux66f2ux7ddaux306eux5e7eux4f55ux66f2ux7387ux637bux7387ux30d5ux30ebux30cdux30bbux30ecux306eux516cux5f0f}{%
\subsection{問題Ⅰ　曲線の幾何（曲率・捻率・フルネ＝セレの公式）}\label{ux554fux984cux2170-ux66f2ux7ddaux306eux5e7eux4f55ux66f2ux7387ux637bux7387ux30d5ux30ebux30cdux30bbux30ecux306eux516cux5f0f}}

次の文章を読んで、【　】に適した式・数値・語句を、\{　\}
からは適切なものを一つ選びその番号を、それぞれ答えよ。また、問1・問2では指示にしたがって答えよ。

\begin{enumerate}
\def\labelenumi{(\arabic{enumi})}
\tightlist
\item
  {\fallbackfont ℝ}³ 内の曲線 c(t) で、すべての t について c′(t) ≠ 0
  を満たすものを正則曲線という。t{\fallbackfont ₀} からの弧長
\end{enumerate}

\[s(t) = \int_{t_0}^{t} |c'(u)|\, du\]

は t の狭義増加関数なので、s
を新しいパラメータにとり直せる。弧長パラメータで表すと、速さは
\textbar dc/ds\textbar{} = 【ア】
となり、曲線の形だけを反映した量を取り出せる。以下 ′ は s
による微分とする。

単位接ベクトルを e{\fallbackfont ₁} = c′(s)
とおく。\textbar e{\fallbackfont ₁}\textbar² = 1 の両辺を s で微分すると
e{\fallbackfont ₁} · e{\fallbackfont ₁}′ = 【イ】
となるので、e{\fallbackfont ₁}′ は e{\fallbackfont ₁}
に直交する。e{\fallbackfont ₁}′
は「進行方向がどれだけ速く向きを変えるか」を表す。

平面曲線では、e{\fallbackfont ₁} を反時計回りに 90° 回転したベクトルを
e{\fallbackfont ₂} とすると e{\fallbackfont ₁}′ = κe{\fallbackfont ₂}
と書け、κ を符号付き曲率という。たとえば半径 r の円 c(s) = (r cos(s/r),
r sin(s/r)) を反時計回りにたどると、e{\fallbackfont ₁} = (−sin(s/r),
cos(s/r))、e{\fallbackfont ₁}′ = −(1/r)(cos(s/r),
sin(s/r))、e{\fallbackfont ₂} = (−cos(s/r), −sin(s/r)) なので κ = 【ウ】
である。一般の曲線でも 1/\textbar κ\textbar{}
は、その点で曲線に最もよく接する円（接触円）の半径に等しい。

\begin{enumerate}
\def\labelenumi{(\arabic{enumi})}
\setcounter{enumi}{1}
\tightlist
\item
  空間曲線で κ = \textbar e{\fallbackfont ₁}′\textbar{} \textgreater{} 0
  のとき、主法線ベクトル e{\fallbackfont ₂} = e{\fallbackfont ₁}′/κ
  と従法線ベクトル e{\fallbackfont ₃} = e{\fallbackfont ₁} ×
  e{\fallbackfont ₂} を定める。正規直交枠 (e{\fallbackfont ₁},
  e{\fallbackfont ₂}, e{\fallbackfont ₃})
  を曲線に沿って動かすと、その変化は次のフルネ＝セレの公式で表される。
\end{enumerate}

\[e_1' = \kappa e_2, \qquad e_2' = -\kappa e_1 + \tau e_3, \qquad e_3' = \text{【エ】}\]

τ
を捩率（れいりつ、捻率とも書く）といい、曲線が接触平面（e{\fallbackfont ₁},
e{\fallbackfont ₂}
の張る平面）から「ねじれて」出ていく速さを表す。係数を並べた 3×3 行列は
【オ】
になっている。これは枠が常に正規直交であることから、e{\fallbackfont ᵢ} ·
e{\fallbackfont ⱼ} = 一定 を微分して得られる性質である。

具体例として常螺旋 c(t) = (a cos t, a sin t, bt)（a \textgreater{}
0）を考える。\textbar c′(t)\textbar{} = 【カ】 は一定なので s は t
の定数倍である。計算すると

\begin{itemize}
\tightlist
\item
  曲率 κ = 【キ】
\item
  捻率 τ = 【ク】
\end{itemize}

となり、どちらも一定である。b = 0 とすると半径 a の円に戻り、κ = 1/a、τ
= 0 となることを確かめておくとよい。逆に、κ と τ を s
の関数として与えると、曲線は空間内の位置と向き（合同変換）を除いて一意に決まる（曲線論の基本定理）。とくに
κ \textgreater{} 0 で τ ≡ 0 である曲線は \{ ケ \}。

\begin{itemize}
\tightlist
\item
  ① ある平面に含まれる
\item
  ② 直線である
\item
  ③ 常螺旋である
\end{itemize}

問1　e{\fallbackfont ₁}′ と e{\fallbackfont ₂}′ の式および
(e{\fallbackfont ₁}, e{\fallbackfont ₂}, e{\fallbackfont ₃})
の正規直交性を用いて、e{\fallbackfont ₃}′ = 【エ】 を導け。

問2　平面曲線 c(t) = (x(t), y(t))
を一般のパラメータで表したとき、符号付き曲率は

\[\kappa = \frac{x'y'' - x''y'}{(x'^2 + y'^2)^{3/2}}\]

で与えられる（′ は t による微分）。この公式を用いて、放物線 y = x²
の原点における曲率と接触円の半径を求めよ。

\hypertarget{ux554fux984cux2171-ux66f2ux9762ux3068ux7b2cux4e00ux57faux672cux5f62ux5f0f}{%
\subsection{問題Ⅱ　曲面と第一基本形式}\label{ux554fux984cux2171-ux66f2ux9762ux3068ux7b2cux4e00ux57faux672cux5f62ux5f0f}}

次の文章を読んで、【　】に適した式・数値を、\{　\}
からは適切なものを一つ選びその番号を、それぞれ答えよ。また、問1・問2では指示にしたがって答えよ。

\begin{enumerate}
\def\labelenumi{(\arabic{enumi})}
\tightlist
\item
  平面領域 D から {\fallbackfont ℝ}³ への滑らかな写像 p(u, v) で、各点で
  p{\fallbackfont ᵤ} と p{\fallbackfont ᵥ}
  が一次独立なものを（正則な）曲面片という。p{\fallbackfont ᵤ},
  p{\fallbackfont ᵥ} が張る平面が接平面であり、単位法線ベクトルは n =
  (p{\fallbackfont ᵤ} × p{\fallbackfont ᵥ})/\textbar p{\fallbackfont ᵤ}
  × p{\fallbackfont ᵥ}\textbar{} である。
\end{enumerate}

曲面上の曲線 c(t) = p(u(t), v(t)) の速度は連鎖律により c′ =
u′p{\fallbackfont ᵤ} + v′p{\fallbackfont ᵥ} なので、その長さの2乗は

\[|c'|^2 = E\,u'^2 + 2F\,u'v' + G\,v'^2, \qquad E = p_u \cdot p_u,\quad F = \text{【イ】},\quad G = p_v \cdot p_v\]

となる。微小な線素で書けば ds² = E du² + 2F du dv + G dv²
であり、これを第一基本形式という。曲面上の曲線の長さは、空間を知らなくても
E, F, G だけから計算できる。

\begin{enumerate}
\def\labelenumi{(\arabic{enumi})}
\setcounter{enumi}{1}
\tightlist
\item
  面積も第一基本形式で表せる。微小な平行四辺形 p{\fallbackfont ᵤ} du,
  p{\fallbackfont ᵥ} dv の面積は \textbar p{\fallbackfont ᵤ} ×
  p{\fallbackfont ᵥ}\textbar{} du dv であり、ベクトルの恒等式 \textbar a
  × b\textbar² = \textbar a\textbar²\textbar b\textbar² − (a · b)²
  を使うと
\end{enumerate}

\[\text{面積} = \iint_D \sqrt{\text{【ロ】}}\; du\, dv\]

となる。半径 R の球面 p(θ, φ) = (R sin θ cos φ, R sin θ sin φ, R cos
θ)（0 \textless{} θ \textless{} π、0 \textless{} φ \textless{}
2π）では、E = 【ハ】、F = 【ニ】、G = 【ホ】 であり、面積は 【ヘ】
となる。また座標曲線（u 一定の曲線と v 一定の曲線）のなす角 ω は cos ω =
F/√(EG) で与えられるので、F = 0 は座標曲線が直交することを意味する。

\begin{enumerate}
\def\labelenumi{(\arabic{enumi})}
\setcounter{enumi}{2}
\tightlist
\item
  2つの曲面の間の対応で、曲線の長さをすべて保つものを等長対応という。適当なパラメータで第一基本形式（E,
  F, G）が一致すれば、2つの曲面は（局所的に）等長である。たとえば平面
  p(u, v) = (u, v, 0) の第一基本形式は du² + dv² である。円柱 q(u, v) =
  (cos u, sin u, v) を計算すると E = 【ト】、F = 0、G = 1
  となって平面と一致する。紙を丸めて円柱を作っても、紙の上の長さが変わらないことの数式表現である。
\end{enumerate}

回転面 p(u, v) = (f(u) cos v, f(u) sin v, g(u))（f \textgreater{}
0）では、E = f′² + g′²、F = 0、G = 【チ】 となる。断面曲線 (f(u), g(u))
を弧長でパラメータ表示しておくと f′² + g′² = 1 なので E = 【リ】
となり、第一基本形式は du² + f(u)² dv²
という簡潔な形になる。この形は問題Ⅳで曲率を計算するときに便利である。

第一基本形式だけで決まる量を、曲面の内在的な量という。次のうち、第一基本形式だけで決まるものは
\{ ヌ \}。

\begin{itemize}
\tightlist
\item
  ① 曲面上の曲線の長さ、交わる角、領域の面積
\item
  ② 単位法線ベクトルの向き
\item
  ③ 空間内での曲面の曲がり方を表す第二基本形式（問題Ⅲ）
\end{itemize}

問1　恒等式 \textbar a × b\textbar² =
\textbar a\textbar²\textbar b\textbar² − (a · b)²
を成分計算などで示し、面積要素が【ロ】の平方根で与えられることを確かめよ。

問2　円錐面（頂点を除く）p(u, v) = (u cos v, u sin v, u)（u
\textgreater{} 0）の第一基本形式を求めよ。さらに、平面の極座標 (r,
{\fallbackfont ϑ}) における第一基本形式 dr² + r² d{\fallbackfont ϑ}²
と比べ、r = √2 u、{\fallbackfont ϑ} = v/√2
と対応させれば円錐面が平面の一部と局所的に等長であることを示せ。

\hypertarget{ux554fux984cux2172-ux7b2cux4e8cux57faux672cux5f62ux5f0fux3068ux30acux30a6ux30b9ux66f2ux7387}{%
\subsection{問題Ⅲ　第二基本形式とガウス曲率}\label{ux554fux984cux2172-ux7b2cux4e8cux57faux672cux5f62ux5f0fux3068ux30acux30a6ux30b9ux66f2ux7387}}

次の文章を読んで、【　】に適した式・数値・語句を、\{　\}
からは適切なものを一つ選びその番号を、それぞれ答えよ。また、問1・問2では指示にしたがって答えよ。

\begin{enumerate}
\def\labelenumi{(\arabic{enumi})}
\tightlist
\item
  第一基本形式は曲面の「内側の長さ」を測るが、曲面が空間内でどう曲がっているかは測らない。それを測るのが、2階微分の法線成分
\end{enumerate}

\[L = p_{uu} \cdot n, \qquad M = p_{uv} \cdot n, \qquad N = p_{vv} \cdot n\]

を係数とする第二基本形式 L du² + 2M du dv + N dv²
である。p{\fallbackfont ᵤ} · n = 0 を u で微分すると L = 【あ】
と書けるので、第二基本形式は「法線ベクトル n
が動く速さ」を表しているともいえる。

曲面上の曲線を弧長パラメータで表すと、加速度 c″ の法線成分
κ{\fallbackfont ₙ} = c″ · n は曲線の速度の方向だけで決まり、

\[\kappa_n = \frac{L\,u'^2 + 2M\,u'v' + N\,v'^2}{E\,u'^2 + 2F\,u'v' + G\,v'^2}\]

となる（′ は任意のパラメータでの微分でよい）。κ{\fallbackfont ₙ}
を法曲率という。

\begin{enumerate}
\def\labelenumi{(\arabic{enumi})}
\setcounter{enumi}{1}
\tightlist
\item
  接平面内の方向を回すと法曲率は変化し、その最大値と最小値を主曲率
  κ{\fallbackfont ₁},
  κ{\fallbackfont ₂}、それを与える方向を主方向という。主曲率は行列
  (第一基本形式の行列){\fallbackfont ⁻}¹(第二基本形式の行列)
  の固有値として求まり、
\end{enumerate}

\[K = \kappa_1 \kappa_2 = \text{【い】}, \qquad H = \frac{\kappa_1 + \kappa_2}{2} = \frac{EN - 2FM + GL}{2(EG - F^2)}\]

をそれぞれガウス曲率、平均曲率という。

半径 R の球面で外向き単位法線 n = p/R をとると、n{\fallbackfont ᵤ} =
p{\fallbackfont ᵤ}/R、n{\fallbackfont ᵥ} = p{\fallbackfont ᵥ}/R
となり、第二基本形式は第一基本形式の −1/R
倍になる。したがってすべての方向で法曲率は等しく、κ{\fallbackfont ₁} =
κ{\fallbackfont ₂} = 【う】、K = 【え】
である。法線の向きを反転すると主曲率の符号は変わるが、K は変わらない。

\begin{enumerate}
\def\labelenumi{(\arabic{enumi})}
\setcounter{enumi}{2}
\tightlist
\item
  曲面が z = f(x, y)、f(0, 0) = 0、∇f(0, 0) = 0
  と表されているとき、原点では E = G = 1、F = 0、n = (0, 0, 1)、L =
  f{\fallbackfont ₓₓ}、M = f{\fallbackfont ₓᵧ}、N = f{\fallbackfont ᵧᵧ}
  となるので、K(0) = 【お】
  である。つまりガウス曲率は、接平面からの「高さ関数」のヘッセ行列の行列式である。
\end{enumerate}

\begin{itemize}
\tightlist
\item
  鞍形 z = xy の原点：K = 【か】
\item
  放物面 z = x² + y² の原点：K = 【き】
\end{itemize}

K \textgreater{} 0 の点を楕円点（近くで曲面が接平面の片側にある）、K
\textless{} 0 の点を 【く】 点（接平面を横切る鞍の形）、K = 0
の点を放物点または平坦点という。半径 r
の円柱面では、主曲率の絶対値は母線方向で 0、円周方向で 1/r なので K =
【け】 である。

なお、球面のように κ{\fallbackfont ₁} = κ{\fallbackfont ₂}
となる点（臍点）以外では主方向はちょうど2つあり、それらは \{ こ \}。

\begin{itemize}
\tightlist
\item
  ① 互いに直交する
\item
  ② 互いに平行である
\item
  ③ 一般には任意の角度をなす
\end{itemize}

問1　鞍形 z = xy の原点における主曲率と主方向を求めよ。

問2　ある点の主曲率を κ{\fallbackfont ₁},
κ{\fallbackfont ₂}、対応する単位主方向を e{\fallbackfont ₁},
e{\fallbackfont ₂} とする。e{\fallbackfont ₁} と角 θ
をなす単位接ベクトル方向の法曲率が

\[\kappa_n(\theta) = \kappa_1 \cos^2\theta + \kappa_2 \sin^2\theta\]

となること（オイラーの公式）を示せ。（ヒント：e{\fallbackfont ₁},
e{\fallbackfont ₂}
を座標軸とする正規直交座標で第一・第二基本形式を表す。）

\hypertarget{ux554fux984cux2173-ux6e2cux5730ux7ddaux9a5aux7570ux306eux5b9aux7406ux30acux30a6ux30b9ux30dcux30cdux306eux5b9aux7406}{%
\subsection{問題Ⅳ　測地線・驚異の定理・ガウス＝ボネの定理}\label{ux554fux984cux2173-ux6e2cux5730ux7ddaux9a5aux7570ux306eux5b9aux7406ux30acux30a6ux30b9ux30dcux30cdux306eux5b9aux7406}}

次の文章を読んで、【　】に適した式・数値・語句を、\{　\}
からは適切なものを一つ選びその番号を、それぞれ答えよ。また、問1・問2では指示にしたがって答えよ。

\begin{enumerate}
\def\labelenumi{(\arabic{enumi})}
\tightlist
\item
  曲面上の曲線
  c(s)（弧長パラメータ）の加速度を、法線方向と、接平面内で速度に直交する方向に分解する。
\end{enumerate}

\[c'' = \kappa_g\, (n \times c') + \kappa_n\, n\]

κ{\fallbackfont ₙ}
は問題Ⅲの法曲率、κg（右辺第1項の係数）を測地的曲率という。曲面上に住む生き物には法線方向は見えないので、κg
こそが「曲面の中で感じる曲がり具合」である。κg ≡ 0、すなわち c″ が常に
【ア】
に平行な曲線を測地線という。曲面上の2点を結ぶ最短の曲線は測地線になる。

半径 R
の球面上の大円を一定の速さでたどると、加速度は球の中心を向くので法線に平行であり、大円は測地線である。一方、赤道以外の緯線では加速度が回転軸を向き、法線と平行にならないので測地線ではない。

\begin{enumerate}
\def\labelenumi{(\arabic{enumi})}
\setcounter{enumi}{1}
\tightlist
\item
  座標 (u¹, u²) = (u, v) と第一基本形式の係数
  g{\fallbackfont ᵢⱼ}（g{\fallbackfont ₁₁} = E、g{\fallbackfont ₁₂} =
  g{\fallbackfont ₂₁} = F、g{\fallbackfont ₂₂} =
  G）を用いると、測地線の方程式は
\end{enumerate}

\[\frac{d^2u^k}{ds^2} + \sum_{i,j} \Gamma^k_{ij} \frac{du^i}{ds} \frac{du^j}{ds} = 0, \qquad \Gamma^k_{ij} = \frac{1}{2} \sum_{l} g^{kl} \left( \partial_i g_{jl} + \partial_j g_{il} - \partial_l g_{ij} \right)\]

と書ける（g{\fallbackfont ᵏˡ} は g{\fallbackfont ᵢⱼ}
の逆行列の成分）。クリストッフェル記号 Γ は 【イ】
とその微分だけで決まるので、測地線も内在的な概念である。

平面の極座標 ds² = dr² + r² d{\fallbackfont ϑ}² で計算すると、0
でないものは Γ{\fallbackfont ʳ}（{\fallbackfont ϑϑ}）=
【ウ】、Γ{\fallbackfont ᶿ}（r{\fallbackfont ϑ}）=
Γ{\fallbackfont ᶿ}（{\fallbackfont ϑ}r）= 【エ】
である。測地線の方程式は r″ − r {\fallbackfont ϑ}′² =
0、{\fallbackfont ϑ}″ + (2/r) r′{\fallbackfont ϑ}′ = 0 となり、第2式は
(r²{\fallbackfont ϑ}′)′ = 0、すなわち r²{\fallbackfont ϑ}′ = 【オ】
を意味する。これは等速直線運動における角運動量保存と同じ式であり、平面の測地線は直線であることが確かめられる。

\begin{enumerate}
\def\labelenumi{(\arabic{enumi})}
\setcounter{enumi}{2}
\tightlist
\item
  ガウスの驚異の定理：ガウス曲率 K
  は、第二基本形式を使って定義したにもかかわらず、第一基本形式とその微分だけで表せる。F
  = 0 の直交座標系では
\end{enumerate}

\[K = -\frac{1}{2\sqrt{EG}} \left[ \frac{\partial}{\partial u}\left( \frac{G_u}{\sqrt{EG}} \right) + \frac{\partial}{\partial v}\left( \frac{E_v}{\sqrt{EG}} \right) \right]\]

である。問題Ⅱの球面（u = θ、v = φ、E = R²、G = R² sin²θ）に適用すると K
= 【カ】 となり、問題Ⅲの結果と一致する。この定理から、K
が異なる曲面どうしは局所的にも等長になりえない。平面は K = 0、球面は K
\textgreater{} 0
なので、地球の表面の一部を長さの比を一切変えずに平面の地図に写すことは
【キ】。一方、円柱面は K = 【ク】
なので、問題Ⅱで見た平面との等長性と矛盾しない。

\begin{enumerate}
\def\labelenumi{(\arabic{enumi})}
\setcounter{enumi}{3}
\tightlist
\item
  3本の測地線で囲まれた曲面上の三角形 T の内角を α, β, γ
  とすると、ガウス＝ボネの定理
\end{enumerate}

\[\iint_T K\, dA = \alpha + \beta + \gamma - \text{【ケ】}\]

が成り立つ。たとえば単位球面上で、北極と赤道上の経度 0°、90°
の2点を頂点とする三角形は3つの角がすべて直角で、面積は球面の 1/8
なので、両辺がともに π/2 になる。この定理により、K \textless{} 0
で一定の曲面（双曲平面）上の測地三角形の内角の和は \{ コ \}。

\begin{itemize}
\tightlist
\item
  ① π より小さい
\item
  ② π に等しい
\item
  ③ π より大きい
\end{itemize}

問1　問題Ⅱの回転面で、断面曲線を弧長パラメータで表したとき、子午線（v =
一定の曲線）が測地線であることを示せ。（ヒント：c″ が p{\fallbackfont ᵤ}
と p{\fallbackfont ᵥ} の両方に直交することを示せばよい。）

問2　単位球面 ds² = dθ² + sin²θ dφ²
のクリストッフェル記号を計算し、測地線に沿って sin²θ · (dφ/ds)
が一定であること（クレローの関係）を導け。

\hypertarget{ux554fux984cux2174-ux53efux5faeux5206ux591aux69d8ux4f53ux3068ux63a5ux7a7aux9593}{%
\subsection{問題Ⅴ　可微分多様体と接空間}\label{ux554fux984cux2174-ux53efux5faeux5206ux591aux69d8ux4f53ux3068ux63a5ux7a7aux9593}}

次の文章を読んで、【　】に適した式・数値・語句を、\{　\}
からは適切なものを一つ選びその番号を、それぞれ答えよ。また、問1・問2では指示にしたがって答えよ。

\begin{enumerate}
\def\labelenumi{(\arabic{enumi})}
\tightlist
\item
  問題Ⅳで見たように、曲面の内在的な幾何は周りの空間がなくても語れる。そこで「各点の近くで
  {\fallbackfont ℝⁿ}
  と同じに見える空間」を直接定義する。ハウスドルフかつ第二可算な位相空間
  M が、開集合 U と同相写像 φ : U → φ(U) ⊂ {\fallbackfont ℝⁿ}
  の組（チャート）で覆われているとき、M を n
  次元位相多様体という。2つのチャートが重なる部分では座標変換 φβ
  {\fallbackfont ∘} φα{\fallbackfont ⁻}¹ が定まり、これがすべて C\^{}∞
  級であるとき M を可微分多様体という。
\end{enumerate}

例として円周 S¹ = \{(x, y) : x² + y² = 1\} を考える。北極 N = (0, 1)
からの立体射影 φN(x, y) = x/(1 − y) は S¹ \textbackslash{} \{N\}
上のチャートであり、南極 S = (0, −1) からの立体射影 φS(x, y) = x/(1 + y)
は S¹ \textbackslash{} \{S\} 上のチャートである。重なり部分で

\[\varphi_N \cdot \varphi_S = \frac{x^2}{1 - y^2} = \text{【イ】}\]

なので、座標変換は t {\fallbackfont ↦} 【ロ】（t ≠ 0）であり、C\^{}∞
級である。S¹ はコンパクトだが {\fallbackfont ℝ}
の空でない開集合はコンパクトでないので、S¹
を1枚のチャートで覆うことはできず、少なくとも 【ハ】 枚が必要である。

\begin{enumerate}
\def\labelenumi{(\arabic{enumi})}
\setcounter{enumi}{1}
\tightlist
\item
  多様体は {\fallbackfont ℝᴺ}
  の中にあるとは限らないので、接ベクトルを「矢印」として定義できない。そこで「方向微分」としての性質を定義に採用する。点
  p での接ベクトルとは、p の近くの滑らかな関数に実数を対応させる線形写像
  v で、ライプニッツ則 v(fg) = v(f) g(p) + f(p) v(g)
  を満たすものである。定数関数 1 については v(1) = v(1·1) = 2v(1) より
  v(1) = 【ニ】 となる。
\end{enumerate}

局所座標 (x¹, \ldots, x{\fallbackfont ⁿ}) のもとで偏微分
∂/∂x{\fallbackfont ⁱ}\textbar{}{\fallbackfont ₚ}
は接ベクトルであり、接ベクトル全体 T{\fallbackfont ₚ}M
はこれらを基底とする 【ホ】 次元のベクトル空間になる。別の座標 (y¹,
\ldots, y{\fallbackfont ⁿ}) をとると、連鎖律から

\[\frac{\partial}{\partial y^j} = \sum_i \frac{\partial x^i}{\partial y^j} \frac{\partial}{\partial x^i}\]

が成り立つ。同じ接ベクトルを v = Σ a{\fallbackfont ⁱ}
∂/∂x{\fallbackfont ⁱ} = Σ b{\fallbackfont ʲ} ∂/∂y{\fallbackfont ʲ}
と2通りに表すと、成分は b{\fallbackfont ʲ} = 【ヘ】
と変換される。物理でいう「反変ベクトル」の変換則である。

\begin{enumerate}
\def\labelenumi{(\arabic{enumi})}
\setcounter{enumi}{2}
\tightlist
\item
  滑らかな写像 F : M → N の点 p での微分 dF{\fallbackfont ₚ} :
  T{\fallbackfont ₚ}M → T\_F(p)N を (dF{\fallbackfont ₚ}v)(g) = v(g
  {\fallbackfont ∘} F) で定める。座標で表すと dF{\fallbackfont ₚ}
  はヤコビ行列である。
\end{enumerate}

F : {\fallbackfont ℝⁿ⁺}¹ → {\fallbackfont ℝ}、F(x) = \textbar x\textbar²
とすると dF{\fallbackfont ₓ}(v) = 2x · v であり、F{\fallbackfont ⁻}¹(1)
上で x ≠ 0 なので 1 は正則値である。正則値定理により球面
S{\fallbackfont ⁿ} = F{\fallbackfont ⁻}¹(1) は n
次元の部分多様体であり、点 p での接空間は
T{\fallbackfont ₚ}S{\fallbackfont ⁿ} = 【ト】 と、dF{\fallbackfont ₚ}
の核として求まる。

関数 f の微分 df は接ベクトル v に v(f)
を対応させる線形形式（余接ベクトル）で、座標の微分は
dx{\fallbackfont ⁱ}(∂/∂x{\fallbackfont ʲ}) = 【チ】
を満たす。また2つのベクトル場 X, Y に対し、{[}X, Y{]}f = X(Yf) − Y(Xf)
で定まる括弧積は、2階微分の項が打ち消し合うので再びベクトル場になる。たとえば
{\fallbackfont ℝ}² で {[}∂/∂x, x ∂/∂y{]} = 【リ】 である。

可微分多様体の定義で座標変換の滑らかさを要求するのは \{ ヌ \}。

\begin{itemize}
\tightlist
\item
  ①
  関数の微分可能性が、どのチャートで調べても同じ結論になるようにするため
\item
  ② 多様体をハウスドルフ空間にするため
\item
  ③ 多様体に距離（計量）を入れるため
\end{itemize}

問1　立体射影 φN の逆写像が t {\fallbackfont ↦} (2t/(t² + 1), (t² −
1)/(t² + 1)) であることを確かめよ（像が S¹ 上にあることと、φN
と合成すると恒等写像になることを示す）。

問2　成分の変換則【ヘ】を導け。また {\fallbackfont ℝ}² で直交座標 (x, y)
と極座標 (r, {\fallbackfont ϑ}) を用いるとき、∂/∂{\fallbackfont ϑ} = −y
∂/∂x + x ∂/∂y を示せ。

\hypertarget{ux554fux984cux2175-ux5faeux5206ux5f62ux5f0fux3068ux30b9ux30c8ux30fcux30afux30b9ux306eux5b9aux7406}{%
\subsection{問題Ⅵ　微分形式とストークスの定理}\label{ux554fux984cux2175-ux5faeux5206ux5f62ux5f0fux3068ux30b9ux30c8ux30fcux30afux30b9ux306eux5b9aux7406}}

次の文章を読んで、【　】に適した式・数値・語句を、\{　\}
からは適切なものを一つ選びその番号を、それぞれ答えよ。また、問1・問2では指示にしたがって答えよ。

\begin{enumerate}
\def\labelenumi{(\arabic{enumi})}
\tightlist
\item
  多様体上で積分を座標によらずに扱うための道具が微分形式である。1次微分形式は
  ω = a{\fallbackfont ₁}dx¹ + \ldots{} +
  a{\fallbackfont ₙ}dx{\fallbackfont ⁿ}
  の形をしており、問題Ⅴの余接ベクトルを点ごとに与えたものである。k
  次形式は dx\^{}i ∧ dx\^{}j ∧ \ldots{} の一次結合で、外積 ∧ は反対称 dx
  ∧ dy = −dy ∧ dx を満たす。したがって dx ∧ dx = 【あ】
  である。{\fallbackfont ℝ}² で計算すると
\end{enumerate}

\[(a\,dx + b\,dy) \wedge (c\,dx + e\,dy) = \text{【い】}\; dx \wedge dy\]

となり、係数に行列式が現れる。2次形式 dx ∧ dy
は「向きのついた面積要素」であり、極座標 x = r cos {\fallbackfont ϑ}、y
= r sin {\fallbackfont ϑ} に変換すると dx ∧ dy = 【う】 dr ∧
d{\fallbackfont ϑ}
となる。重積分の変数変換でヤコビアンが現れる理由は、この計算に集約されている。

\begin{enumerate}
\def\labelenumi{(\arabic{enumi})}
\setcounter{enumi}{1}
\tightlist
\item
  外微分 d は k 次形式を (k + 1) 次形式に写す。関数（0次形式）には df =
  Σ (∂f/∂x{\fallbackfont ⁱ}) dx{\fallbackfont ⁱ}、一般には d(a dx\^{}I)
  = da ∧ dx\^{}I
  と定める。偏微分の順序交換が可能なことと外積の反対称性から、つねに d
  {\fallbackfont ∘} d = 【え】 となる。
\end{enumerate}

{\fallbackfont ℝ}³
では1次形式・2次形式をベクトル場と同一視でき、外微分は次のように対応する。

\begin{longtable}[]{@{}lll@{}}
\toprule\noalign{}
形式 & d の作用 & ベクトル解析での対応 \\
\midrule\noalign{}
\endhead
\bottomrule\noalign{}
\endlastfoot
0次形式 f & df & grad f \\
1次形式 & dω & 【お】 \\
2次形式 & dω & div \\
\end{longtable}

d {\fallbackfont ∘} d = 【え】 は、ベクトル解析の公式 rot(grad f) = 0 と
div(rot A) = 0 をまとめて表している。

\begin{enumerate}
\def\labelenumi{(\arabic{enumi})}
\setcounter{enumi}{2}
\tightlist
\item
  ストークスの定理：境界 ∂M をもつ向き付けられた k 次元多様体 M と (k −
  1) 次形式 ω に対し
\end{enumerate}

\[\int_M d\omega = \int_{\partial M} \omega\]

が成り立つ。k = 1 なら微積分の基本定理そのものである。平面領域 D と ω =
P dx + Q dy の場合、dω = 【か】 dx ∧ dy
となり、グリーンの定理が得られる。とくに ω = (x dy − y dx)/2 とすると dω
= dx ∧ dy なので、D の面積は境界の線積分 (1/2)∮(x dy − y dx)
で求まる。半径 a の円板で x = a cos t、y = a sin t とおくと被積分関数は
a² dt となり、面積は 【き】
である。ガウスの発散定理や古典的なストークスの定理も、この1つの式の特別な場合である。

\begin{enumerate}
\def\labelenumi{(\arabic{enumi})}
\setcounter{enumi}{3}
\tightlist
\item
  dω = 0 となる形式を閉形式、ω = dη と書ける形式を完全形式という。d
  {\fallbackfont ∘} d = 0
  より完全形式は閉形式である。では逆はどうか。{\fallbackfont ℝ}²
  \textbackslash{} \{0\} 上の1次形式
\end{enumerate}

\[\omega = \frac{-y\,dx + x\,dy}{x^2 + y^2}\]

を考えると、計算により dω = 【く】
である。しかし単位円を反時計回りに一周する積分は ∮ω = 【け】 であり 0
でない。もし ω = dη と書けたなら、ストークスの定理（k =
1）より閉曲線上の積分は 0 になるはずなので、ω
は完全形式ではない。原点という「穴」があるために閉形式が完全形式にならないのであり、閉形式と完全形式のずれが空間の位相を検出する（ド・ラーム・コホモロジー）。一方、{\fallbackfont ℝⁿ}
全体や星形の領域では、閉形式は \{ こ \}。

\begin{itemize}
\tightlist
\item
  ① 必ず完全形式である（ポアンカレの補題）
\item
  ② 必ずしも完全形式ではない
\item
  ③ 必ず 0 である
\end{itemize}

問1　{\fallbackfont ℝ}² 上の滑らかな関数 f について、d(df) = 0
を直接計算で確かめよ。

問2　(4) の ω は、局所的には極座標の角度 {\fallbackfont ϑ} を用いて ω =
d{\fallbackfont ϑ} と書けることを示せ。また、それにもかかわらず ω が
{\fallbackfont ℝ}² \textbackslash{} \{0\}
全体で完全形式にならない理由を、{\fallbackfont ϑ} の多価性と ∮ω = 【け】
を結びつけて説明せよ。

\hypertarget{ux89e3ux7b54ux3068ux89e3ux8aac}{%
\subsection{解答と解説}\label{ux89e3ux7b54ux3068ux89e3ux8aac}}

空欄は表で、問は答案例で示す。答案例は要点を押さえた1つの書き方であり、これに限らない。

\hypertarget{ux554fux984cux2170}{%
\subsubsection{問題Ⅰ}\label{ux554fux984cux2170}}

\begin{longtable}[]{@{}lll@{}}
\toprule\noalign{}
空欄 & 解答 & ポイント \\
\midrule\noalign{}
\endhead
\bottomrule\noalign{}
\endlastfoot
ア & 1 & \\
イ & 0 & \\
ウ & 1/r & 反時計回りなら正 \\
エ & −τe{\fallbackfont ₂} & 問1 \\
オ & 交代行列（反対称行列） & e{\fallbackfont ᵢ} · e{\fallbackfont ⱼ} =
δ{\fallbackfont ᵢⱼ} を微分 \\
カ & √(a² + b²) & \\
キ & a/(a² + b²) & \\
ク & b/(a² + b²) & b の符号で右巻き・左巻きが決まる \\
ケ & ① & 従法線 e{\fallbackfont ₃} が一定 \\
\end{longtable}

常螺旋の計算：w = √(a² + b²)、s = wt とおくと e{\fallbackfont ₁} = (−a
sin t, a cos t, b)/w、de{\fallbackfont ₁}/ds = (−a cos t, −a sin t,
0)/w² より κ = a/w²、e{\fallbackfont ₂} = (−cos t, −sin t,
0)。e{\fallbackfont ₃} = e{\fallbackfont ₁} × e{\fallbackfont ₂} = (b
sin t, −b cos t, a)/w、de{\fallbackfont ₃}/ds = (b cos t, b sin t, 0)/w²
= −(b/w²)e{\fallbackfont ₂} より τ = b/w²。

問1　e{\fallbackfont ₃} は単位ベクトルなので e{\fallbackfont ₃}′ ⊥
e{\fallbackfont ₃}。また e{\fallbackfont ₃} · e{\fallbackfont ₁} = 0
を微分すると e{\fallbackfont ₃}′ · e{\fallbackfont ₁} =
−e{\fallbackfont ₃} · e{\fallbackfont ₁}′ = −κ e{\fallbackfont ₃} ·
e{\fallbackfont ₂} = 0。よって e{\fallbackfont ₃}′ は e{\fallbackfont ₂}
方向で、係数は e{\fallbackfont ₃}′ · e{\fallbackfont ₂} =
−e{\fallbackfont ₃} · e{\fallbackfont ₂}′ = −τ。したがって
e{\fallbackfont ₃}′ = −τe{\fallbackfont ₂}。

問2　x = t、y = t² とおくと x′ = 1、x″ = 0、y′ = 2t、y″ = 2。t = 0 で κ
= (1·2 − 0)/1 = 2。接触円の半径は 1/2。

\hypertarget{ux554fux984cux2171}{%
\subsubsection{問題Ⅱ}\label{ux554fux984cux2171}}

\begin{longtable}[]{@{}lll@{}}
\toprule\noalign{}
空欄 & 解答 & ポイント \\
\midrule\noalign{}
\endhead
\bottomrule\noalign{}
\endlastfoot
イ & p{\fallbackfont ᵤ} · p{\fallbackfont ᵥ} & \\
ロ & EG − F² & \\
ハ & R² & \\
ニ & 0 & 経線と緯線は直交 \\
ホ & R² sin²θ & 緯線の半径は R sin θ \\
ヘ & 4πR² & ∫∫ R² sin θ dθ dφ \\
ト & 1 & \\
チ & f(u)² & \\
リ & 1 & \\
ヌ & ① & \\
\end{longtable}

問1　a と b のなす角を θ とすると \textbar a × b\textbar{} =
\textbar a\textbar\textbar b\textbar{} sin θ、a · b =
\textbar a\textbar\textbar b\textbar{} cos θ なので、sin²θ + cos²θ = 1
より恒等式が従う。a = p{\fallbackfont ᵤ}、b = p{\fallbackfont ᵥ}
とおけば \textbar p{\fallbackfont ᵤ} × p{\fallbackfont ᵥ}\textbar² = EG
− F²。

問2　p{\fallbackfont ᵤ} = (cos v, sin v, 1)、p{\fallbackfont ᵥ} = (−u
sin v, u cos v, 0) より E = 2、F = 0、G = u²。ds² = 2 du² + u² dv²。r =
√2 u、{\fallbackfont ϑ} = v/√2 とおくと dr² = 2 du²、r²
d{\fallbackfont ϑ}² = 2u² · dv²/2 = u² dv² となり一致する。v が 0〜2π
を動くとき {\fallbackfont ϑ} は 0〜√2π を動くので、円錐面は中心角 √2π
の扇形を丸めたものと等長である。

\hypertarget{ux554fux984cux2172}{%
\subsubsection{問題Ⅲ}\label{ux554fux984cux2172}}

\begin{longtable}[]{@{}lll@{}}
\toprule\noalign{}
空欄 & 解答 & ポイント \\
\midrule\noalign{}
\endhead
\bottomrule\noalign{}
\endlastfoot
あ & −p{\fallbackfont ᵤ} · n{\fallbackfont ᵤ} & \\
い & (LN − M²)/(EG − F²) & 行列式の比 \\
う & −1/R & 外向き法線の場合 \\
え & 1/R² & \\
お & f{\fallbackfont ₓₓ} f{\fallbackfont ᵧᵧ} − f{\fallbackfont ₓᵧ}² & \\
か & −1 & \\
き & 4 & \\
く & 双曲 & \\
け & 0 & \\
こ & ① & 第一基本形式に関して対称な行列の固有ベクトル \\
\end{longtable}

問1　原点で第一基本形式は単位行列、L = 0、M = 1、N = 0 なので、行列
{[}{[}0, 1{]}, {[}1, 0{]}{]} の固有値を求めて κ{\fallbackfont ₁} =
1、κ{\fallbackfont ₂} = −1（n = (0, 0, 1)）。主方向は接平面（xy
平面）内の (1, 1)/√2 と (1, −1)/√2 で、直交している。

問2　e{\fallbackfont ₁}, e{\fallbackfont ₂}
を基底とする接平面の正規直交座標では、第一基本形式は単位行列、第二基本形式は主方向で対角化されて
diag(κ{\fallbackfont ₁}, κ{\fallbackfont ₂}) となる。方向 (cos θ, sin θ)
を法曲率の式に代入すると、分母は 1、分子は κ{\fallbackfont ₁}cos²θ +
κ{\fallbackfont ₂}sin²θ。

\hypertarget{ux554fux984cux2173}{%
\subsubsection{問題Ⅳ}\label{ux554fux984cux2173}}

\begin{longtable}[]{@{}lll@{}}
\toprule\noalign{}
空欄 & 解答 & ポイント \\
\midrule\noalign{}
\endhead
\bottomrule\noalign{}
\endlastfoot
ア & 法線ベクトル n & \\
イ & 第一基本形式（g{\fallbackfont ᵢⱼ}） & \\
ウ & −r & −(1/2) ∂(r²)/∂r \\
エ & 1/r & (1/2)(1/r²) ∂(r²)/∂r \\
オ & 一定（定数） & \\
カ & 1/R² & G/√(EG) = 2 cos θ を微分 \\
キ & できない & どんな地図にも歪みが出る \\
ク & 0 & \\
ケ & π & 平面なら内角の和 = π \\
コ & ① & 左辺が負 \\
\end{longtable}

問1　子午線 c(u) = (f(u) cos v{\fallbackfont ₀}, f(u) sin
v{\fallbackfont ₀}, g(u)) は弧長パラメータなので c″ ⊥ c′ =
p{\fallbackfont ᵤ}。また c″ = (f″ cos v{\fallbackfont ₀}, f″ sin
v{\fallbackfont ₀}, g″) と p{\fallbackfont ᵥ} = (−f sin
v{\fallbackfont ₀}, f cos v{\fallbackfont ₀}, 0) の内積は 0。よって c″
は接平面に直交し n に平行なので、子午線は測地線である。

問2　g{\fallbackfont ₁₁} = 1、g{\fallbackfont ₂₂} =
sin²θ、g{\fallbackfont ₁₂} = 0 より、0 でないクリストッフェル記号は
Γ\^{}θ(φφ) = −sin θ cos θ、Γ\^{}φ(θφ) = Γ\^{}φ(φθ) = cos θ/sin θ。φ
の測地線方程式は φ″ + 2(cos θ/sin θ) θ′φ′ = 0。一方 (sin²θ φ′)′ = sin²θ
φ″ + 2 sin θ cos θ θ′φ′ = sin²θ (φ″ + 2(cos θ/sin θ) θ′φ′) = 0。よって
sin²θ φ′
は一定。回転面一般で成り立つクレローの関係の特別な場合で、問題Ⅳ(2)
の角運動量保存と同じ構造である。

\hypertarget{ux554fux984cux2174}{%
\subsubsection{問題Ⅴ}\label{ux554fux984cux2174}}

\begin{longtable}[]{@{}lll@{}}
\toprule\noalign{}
空欄 & 解答 & ポイント \\
\midrule\noalign{}
\endhead
\bottomrule\noalign{}
\endlastfoot
イ & 1 & x² = 1 − y² \\
ロ & 1/t & \\
ハ & 2 & \\
ニ & 0 & \\
ホ & n & \\
ヘ & Σ{\fallbackfont ᵢ} (∂y{\fallbackfont ʲ}/∂x{\fallbackfont ⁱ})
a{\fallbackfont ⁱ} & \\
ト & \{v ∈ {\fallbackfont ℝⁿ⁺}¹ : p · v = 0\} &
球面の接平面は半径に直交 \\
チ & δ{\fallbackfont ⁱⱼ}（i = j のとき 1、それ以外 0） & \\
リ & ∂/∂y & \\
ヌ & ① & \\
\end{longtable}

問1　x = 2t/(t² + 1)、y = (t² − 1)/(t² + 1) とおくと x² + y² = (4t² +
t{\fallbackfont ⁴} − 2t² + 1)/(t² + 1)² = 1 で、y \textless{} 1 なので
S¹ \textbackslash{} \{N\} 上にある。1 − y = 2/(t² + 1) より x/(1 − y) =
t となり、φN と合成すると恒等写像。

問2　v を座標関数 y{\fallbackfont ʲ} に作用させる。v = Σ
a{\fallbackfont ⁱ} ∂/∂x{\fallbackfont ⁱ} からは v(y{\fallbackfont ʲ}) =
Σ a{\fallbackfont ⁱ} ∂y{\fallbackfont ʲ}/∂x{\fallbackfont ⁱ}、v = Σ
b{\fallbackfont ᵏ} ∂/∂y{\fallbackfont ᵏ} からは v(y{\fallbackfont ʲ}) =
b{\fallbackfont ʲ}。両者を比べて変換則を得る。極座標では x = r cos
{\fallbackfont ϑ}、y = r sin {\fallbackfont ϑ} より ∂/∂{\fallbackfont ϑ}
= (∂x/∂{\fallbackfont ϑ}) ∂/∂x + (∂y/∂{\fallbackfont ϑ}) ∂/∂y = −r sin
{\fallbackfont ϑ} ∂/∂x + r cos {\fallbackfont ϑ} ∂/∂y = −y ∂/∂x + x
∂/∂y。

\hypertarget{ux554fux984cux2175}{%
\subsubsection{問題Ⅵ}\label{ux554fux984cux2175}}

\begin{longtable}[]{@{}lll@{}}
\toprule\noalign{}
空欄 & 解答 & ポイント \\
\midrule\noalign{}
\endhead
\bottomrule\noalign{}
\endlastfoot
あ & 0 & \\
い & ae − bc & 2×2 行列式 \\
う & r & 極座標のヤコビアン \\
え & 0 & \\
お & rot（curl） & \\
か & ∂Q/∂x − ∂P/∂y & \\
き & πa² & \\
く & 0 & \\
け & 2π & \\
こ & ① & \\
\end{longtable}

\begin{enumerate}
\def\labelenumi{(\arabic{enumi})}
\item
  の極座標の計算：dx = cos {\fallbackfont ϑ} dr − r sin
  {\fallbackfont ϑ} d{\fallbackfont ϑ}、dy = sin {\fallbackfont ϑ} dr +
  r cos {\fallbackfont ϑ} d{\fallbackfont ϑ} より dx ∧ dy = r
  cos²{\fallbackfont ϑ} dr ∧ d{\fallbackfont ϑ} − r
  sin²{\fallbackfont ϑ} d{\fallbackfont ϑ} ∧ dr = r dr ∧
  d{\fallbackfont ϑ}。
\item
  の dω = 0 の計算：P = −y/(x² + y²)、Q = x/(x² + y²) とすると ∂Q/∂x =
  ∂P/∂y = (y² − x²)/(x² + y²)²。
\end{enumerate}

問1　df = f{\fallbackfont ₓ} dx + f{\fallbackfont ᵧ} dy より d(df) =
f{\fallbackfont ₓᵧ} dy ∧ dx + f{\fallbackfont ᵧₓ} dx ∧ dy =
(f{\fallbackfont ᵧₓ} − f{\fallbackfont ₓᵧ}) dx ∧ dy =
0（偏微分の順序交換）。

問2　x \textgreater{} 0 の範囲などで {\fallbackfont ϑ} = arctan(y/x)
とすると、d{\fallbackfont ϑ} = (x dy − y dx)/(x² + y²) = ω。他の範囲でも
{\fallbackfont ϑ} に定数を加えた局所的な表示で同じ式になる。もし
{\fallbackfont ℝ}² \textbackslash{} \{0\} 全体で ω = dη なら η
は局所的に {\fallbackfont ϑ} + 定数に等しく、単位円を一周しても η
は元の値に戻るので ∮ω = 0 となるはずである。しかし実際は 2π
であり、{\fallbackfont ϑ} が一周で 2π
増える（1価関数として定まらない）ことが原因である。

\end{document}
